\subsection{普遍可测函数演算}

在本节中，我们所考虑的希尔伯特空间都是复的。

设 u 是希尔伯特空间 E 上的正规部分算子。我们用 \(\mu_{x,y}\) （相应地，\(\mu_{y}\) ）表示 \((x,y)\in\mathrm{E}\times\mathrm{E}\) （相应地，y）关于 u 的谱测度。

设 \(y \in \mathrm{E}\) 。映射 \(\mathcal{K}(\mathrm{Sp}(u)) \to \mathrm{E}\) 由 \(f \mapsto f(u)(y)\) 所定义，它满足

\[
\| f (u) (y) \| ^ {2} = \int_ {\mathrm{Sp} (u)} | f | ^ {2} \mu_ {y} = \| f \| _ {\mathcal {L} ^ {2} (\mathrm{Sp} (u), \mu_ {y})} ^ {2}.
\]

因此，存在唯一的等距线性映射 \(ev_{y}\) ，它从空间 \(\mathrm{L}^{2}(\mathrm{Sp}(u),\mu_{y})\) 映入 E，使得 \(\mathrm{ev}_{y}(\widetilde{f})=f(u)(y)\) ，这里 \(\widetilde{f}\) 是函数 \(f\in\mathcal{K}(\mathrm{Sp}(u))\) 的类。

设 \(f\) 是定义在 \(u\) 的谱上的普遍可测函数。我们用 \(\mathrm{D}_{f}\) 表示这样的元素 \(y \in \mathrm{E}\) 所成的集：\(f\) 属于 \(\mathcal{L}^{2}(\mathrm{Sp}(u), \mu_{y})\) 。

命题 3. --- 设 f 是定义在 u 的谱上的普遍可测函数。集合 D \(_{f}\) 是 E 的稠密向量子空间。映射 y ↦ ev \(_{y}\) (f) 是 E 上的正规部分算子，其定义域为 D \(_{f}\) ，记作 f(u)。

对每个 \(x \in \mathrm{E}\) 及每个 \(y \in \mathrm{D}_f\) ，我们有 \(f \in \mathcal{L}^1(\mathrm{Sp}(u), \mu_{x,y})\) ，并且

\[
\langle x \mid f (u) y \rangle = \int_ {\operatorname{Sp} (u)} f \mu_ {x, y}.\tag{6}
\]

我们将证明下列更精确的陈述：

命题 4. --- 设 f 是 u 的谱上的普遍可测函数。

\begin{enumerate}
\def\labelenumi{\alph{enumi})}
\item
  集合 \(D_{f}\) 是 E 的稠密向量子空间。映射 \(f(u): y \mapsto \operatorname{ev}_{y}(f)\) 从 \(D_{f}\) 到 E 中是线性的，且当 f = 1 时与 \(1_{E}\) 重合；
\item
  设 \(g \in \mathcal{L}_{\mathrm{u}}(\mathrm{Sp}(u))\) 。对每个 \(y \in \mathrm{D}_f\) 及每个 \(x \in \mathrm{D}_g\) ，我们有 \(fg \in \mathcal{L}^1(\mathrm{Sp}(u), \mu_{x,y})\) ，并且
\end{enumerate}

\[
\langle g (u) x \mid f (u) y \rangle = \int_ {\mathrm{Sp} (u)} \overline {{g}} f \mu_ {x, y}.\tag{7}
\]

\begin{enumerate}
\def\labelenumi{\alph{enumi})}
\setcounter{enumi}{2}
\tightlist
\item
  假设 \(E = L^{2}(X, \mu)\) ，其中 X 是局部紧拓扑空间，\(\mu\) 是 X 上的正测度；又假设 \(u = m_{h}\) ，其中 \(h: X \to \mathbf C\) 是 \(\mu\)-可测的。我们有 \(f(m_{h}) = m_{f \circ h}\) 。
\end{enumerate}

由 IV, p. 266 的定理 1，我们可以假设处于断言 \(c\) ) 的情形，即 \(\mathrm{E} = \mathrm{L}^2 (\mathrm{X},\mu)\) 且 \(u = m_h\) ，其中 X 是局部紧拓扑空间，\(\mu\) 是 X 上的正测度，而 \(h\colon \mathrm{X}\to \mathbf{C}\) 是 \(\mu\)-可测的。我们将用 \(\widetilde{\varphi}\) 表示 \(\mathrm{L}^2 (\mathrm{X},\mu)\) 中函数 \(\varphi \in \mathcal{L}^2 (\mathrm{X},\mu)\) 的类。

设 \(\mathrm{S} = \mathrm{Sp}(m_h)\) 。我们用 \(\mu_{\widetilde{\varphi}_1, \widetilde{\varphi}_2}\) （相应地，\(\mu_{\widetilde{\varphi}}\) ）表示 \((\widetilde{\varphi}_1, \widetilde{\varphi}_2)\) （相应地，\(\widetilde{\varphi}\) ）关于 \(m_h\) 的谱测度，对所有 \((\varphi_1, \varphi_2) \in \mathcal{L}^2(\mathrm{X}, \mu)^2\) （相应地，每个 \(\varphi \in \mathcal{L}^2(\mathrm{X}, \mu)\) ）。

设 \(\varphi\in\mathcal{L}^{2}(X,\mu)\) 。谱测度 \(\mu_{\widetilde{\varphi}}\) 等于 S 上的像测度 \(h(|\varphi|^{2}\cdot\mu)\) （IV, p. 269 的引理 4）。由于函数 \(f\) 是 \(\mu_{\varphi}\)-可测的，我们有 \(\varphi\in D_{f}\) 当且仅当积分

\[
\int_ {\mathrm{S}} ^ {*} | f | ^ {2} d \mu_ {\widetilde {\varphi}} = \int_ {\mathrm{S}} ^ {*} | f | ^ {2} h (| \varphi | ^ {2} d \mu) = \int_ {\mathrm{X}} ^ {*} | f \circ h | ^ {2} | \varphi | ^ {2} d \mu
\]

都是有限的（INT, V, p. 70, § 6, 第 2 目, 命题 2）。这表明 \(D_{f}\) 是乘法算子 \(m_{f\circ h}\) 的定义域，而该定义域是 E 的稠密子空间（IV, p. 232 的命题 5）。

\(\mathrm{ev}_{\widetilde{\varphi}}\) 在函数类 \(g \in \mathcal{K}(\mathrm{S})\) 上的限制是把类 \(\widetilde{g}\) 对应到元素 \(g(m_h)(\widetilde{\varphi}) = m_{g \circ h}(\widetilde{\varphi})\) 的映射（IV, p. 269 的引理 4）。因此，映射 \(\mathrm{ev}_{\widetilde{\varphi}}\) 与从 \(\mathrm{L}^2(\mathrm{S}, \mu_{\widetilde{\varphi}})\) 到 \(\mathrm{L}^2(\mathrm{X}, \mu)\) 中的等距映射重合，后者是由映射 \(g \mapsto (g \circ h) \cdot \varphi\) （从 \(\mathcal{L}^2(\mathrm{S}, \mu_{\widetilde{\varphi}})\) 到 \(\mathcal{L}^2(\mathrm{X}, \mu)\) 中）取商得到的。

特别地，我们因此有 \(\mathrm{ev}_{\widetilde{\varphi}}(f)=m_{f\circ h}(\widetilde{\varphi})\) ，于是 \(f(m_{h})\) 作为从 \(D_{f}\) 到 E 中的映射与部分算子 \(m_{f\circ h}\) 重合。这就证明了断言 c)；若 f=1，则得到 \(\mathrm{ev}_{\widetilde{\varphi}}(1)=\widetilde{\varphi}\) ，这也就完成了 a) 的证明。

我们来证明断言 \(b\) )。设 \(\varphi_{1}\) 与 \(\varphi_{2}\) 是使 \(\widetilde{\varphi}_{1} \in \mathrm{D}_{f}\) 且 \(\widetilde{\varphi}_{2} \in \mathrm{D}_{g}\) 的函数。由（INT, V, loc. cit.）可得

\[
\begin{array}{l} \int_ {\mathrm{S}} ^ {*} | f g |   | \mu_ {\widetilde {\varphi_ {1}}, \widetilde {\varphi_ {2}}} | = \int_ {\mathrm{X}} ^ {*} | (f \circ h) (g \circ h) \varphi_ {1} \varphi_ {2} | d \mu \\ \qquad \leqslant \| (f \circ h) \varphi_ {1} \| \| (g \circ h) \varphi_ {2} \| \end{array}
\]

它是有限的，因此 \(fg \in \mathcal{L}^1(\mathrm{S}, \mu_{\widetilde{\varphi}_1, \widetilde{\varphi}_2})\) 。此外，这时我们有

\[
\begin{array}{c} \langle g (m _ {h}) (\widetilde {\varphi} _ {2})   |   f (m _ {h}) (\widetilde {\varphi} _ {1}) \rangle = \int_ {\mathrm{X}} \overline {{(g \circ h)   \varphi_ {2}}}    (f \circ h)   \varphi_ {1}    d \mu \\ = \int_ {\mathrm{S}} \overline {{g}}   f    d \mu_ {\widetilde {\varphi} _ {1}, \widetilde {\varphi} _ {2}}, \end{array}
\]

证明完毕。

定义 5. --- 映射 \(f \mapsto f(u)\) 由 IV, p. 271 的命题 3 所定义，它从 \(\mathcal{L}_{\mathrm{u}}(\mathrm{Sp}(u))\) 映到 E 上正规部分算子所成的集，称为与 \(u\) 相伴的普遍可测函数演算映射。

更一般地，设 T 是一个集合，并设 \(g\colon\operatorname{Sp}(u)\times\operatorname{T}\to\mathbf{C}\) 是一个映射；设 \(t\in T\) 使得函数 \(g_{t}\colon z\mapsto g(z,t)\) 在 \(\operatorname{Sp}(u)\) 上普遍可测。这时我们记 \(g(u,t)=g_{t}(u)\) 。

公式 (7) 特别地蕴含

\[
\| f (u) (y) \| ^ {2} = \int_ {\operatorname{Sp} (u)} | f | ^ {2} \mu_ {y}\tag{8}
\]

对所有 \(f \in \mathcal{L}_{\mathrm{u}}(\mathrm{Sp}(u))\) 及每个 \(y \in \mathrm{D}_f\) 成立。由此，取 \(f = 1\) ，我们推出 \(\mu_y\) 是总质量为 \(\| y\|^2\) 的正测度。

若 \(u\) 是希尔伯特空间 E 上的正规部分算子且 \(f \in \mathcal{K}(\mathrm{Sp}(u))\) ，则 \(\mathrm{D}_f = \mathrm{E}\) ，且部分算子 \(f(u)\) 这时是连续的，因为

\[
\| f (u) y \| \leqslant \| f \| _ {\infty} \| y \|
\]

对所有 \(y \in E\) 成立（依据 (8)）。自同态 \(f(u)\) 与 IV, p. 268 的定义 3 中的那个自同态重合。

推论 1. --- a) 对每个 \(f \in \mathcal{L}_{\mathrm{u}}(\mathrm{Sp}(u))\) ，部分算子 \(f(u)\) 是正规的，且 \(f(u)^{*} = \overline{f}(u)\) 。此外，\(f(u)\) 在 \(f \geqslant 0\) 时是正的，在 \(f\) 取实值时是自伴的；

\begin{enumerate}
\def\labelenumi{\alph{enumi})}
\setcounter{enumi}{1}
\item
  设 \(k \in \mathbf{N}\) ，令 \(f(z) = z^{k}\) ，\(z \in \operatorname{Sp}(u)\) 。我们有 \(f(u) = u^{k}\) ；
\item
  设 \(\lambda \in \mathbf{C} - \mathrm{Sp}(u)\) ，令 \(f(z) = (\lambda - z)^{-1}\) ，\(z \in \mathrm{Sp}(u)\) 。我们有 \(f(u) = \mathrm{R}(u, \lambda)\) ；
\item
  我们有 \(\beta(u) = b(u)\) ；
\item
  设 \(f \in \mathcal{L}_{\mathrm{u}}(\mathrm{Sp}(u))\) 与 \(g \in \mathcal{L}_{\mathrm{u}}(\mathrm{Sp}(u))\) 满足 \(|g| \leqslant 1 + |f|\) 。\(g(u)\) 的定义域是 \(f(u)\) 的核。
\end{enumerate}

鉴于谱定理（IV, p. 266 的定理 1），这由前面的命题分别结合以下各项而得：

\begin{enumerate}
\def\labelenumi{\alph{enumi})}
\item
  IV, p. 266 的引理 2, a)、IV, p. 253 的命题 23 及其推论；
\item
  IV, p. 255 的命题 24, b)；
\item
  IV, p. 252 的命题 22, b)；
\item
  IV, p. 266 的引理 2, a)；
\item
  IV, p. 232 的命题 6, b)。
\end{enumerate}

注. --- 对 \(f = \operatorname{Id}_{\operatorname{Sp}(u)}\) ，我们有 \(f(u) = u\) （断言 b) 取 \(k = 1\) ）。因此，u 的定义域与这样的 \(x \in E\) 所成的集重合：\(\operatorname{Sp}(u)\) 的恒等映射属于 \(\mathcal{L}^{2}(\operatorname{Sp}(u), \mu_{x})\) ；它特别包含这样的元素 \(x \in E\) ：测度 \(\mu_{x}\) 具有紧支撑。

下面的推论推广了 IV, p. 247 的命题 16 的推论以及 IV, p. 248 的命题 17。

推论 2. --- 设 u 是希尔伯特空间 E 上的正规部分算子。假设 E 非零。

\begin{enumerate}
\def\labelenumi{\alph{enumi})}
\item
  u 的谱非空；
\item
  对每个 \(\lambda \in \mathbf{C} - \mathrm{Sp}(u)\) ，预解式 \(\mathrm{R}(u, \lambda)\) 的范数等于 \(1/\delta\) ，其中 \(\delta > 0\) 是在 \(\mathbf{C}\) 中从 \(\lambda\) 到 \(u\) 的谱的距离。
\item
  对每个 \(\varepsilon > 0\) ，\(\varepsilon\)-伪谱 \(\mathrm{PSp}_{\varepsilon}(u)\) 是由 \(\lambda \in \mathbf{C}\) （距离 \(< \varepsilon\) ，从 \(\mathrm{Sp}(u)\) 量起）所成的集。
\end{enumerate}

假设 \(u\) 的谱为空。那么 \(u\) 是单射的，且自同态 \(u^{-1} = -\mathrm{R}(u,0)\) 是正规的（推论 1, c) 与 a)）。由于 \(\mathrm{Sp}(u^{-1}) \subset \{0\}\) （命题 15, a)），我们有 \(\mathrm{Sp}(u^{-1}) = \{0\}\) （I, p. 26, 推论 1）。由于 \(u^{-1}\) 是正规的，这就蕴含 \(u^{-1} = 0\) （I, p. 110, 例 1），而这是矛盾，因为 E 不为零。

为证明 \(b\) )，可以假设 \(u\) 是乘法算子 \(m_g\) ，作用于 \(\mathrm{L}^2 (\mathrm{X},\mu)\) ，其中 \(g\) 是赋以正测度 \(\mu\) 的局部紧拓扑空间 X 上的连续函数（IV, p. 266 的定理 1）。设 \(\lambda \in \mathbf{C} - \mathrm{Sp}(u)\) ，并设 \(\delta >0\) 是 \(\lambda\) 到 \(u\) 的谱的距离。为证明 \(\| \mathrm{R}(u,\lambda)\| = \delta^{-1}\) ，我们归结到 \(\lambda = 0\) 的情形，办法是把 \(u\) 换成 \(u - \lambda 1_{\mathrm{E}}\) 。于是实数 \(\delta\) 是 0 到 \(m_g\) 的谱的距离。由于后者与以 \(\mu\) 为测度的 \(g\) 的本质像重合，结果是 IV, p. 182 的引理 3 的推论。

最后，断言 \(c\) ) 由 \(b\) ) 及 \(\mathrm{PSp}_{\varepsilon}(u)\) 的定义（IV, p. 250, 定义 9）得出。

推论 3. --- 设 u 是希尔伯特空间 E 上的正规部分算子。设 \(f \in \mathcal{L}_{\mathrm{u}}(\mathrm{Sp}(u))\) 。对 E 中所有 x 与 y，\((x, y)\) 关于 \(f(u)\) 的谱测度是像测度 \(f(\mu)\) ，其中 \(\mu\) 是 \((x, y)\) 关于 u 的谱测度。

鉴于谱定理（IV, p. 266 的定理 1），可以假设 \(u\) 是乘法算子 \(m_g\) （作用于 \(\mathrm{L}^2 (\mathrm{X},\mu)\) ），其中 X 是局部紧拓扑空间，\(\mu\) 是 X 上的正测度，且 \(g\in \mathcal{C}(\mathrm{X})\) 。设 \(f_{1}\) 与 \(f_{2}\) 属于 \(\mathcal{L}^2 (\mathrm{X},\mu)\) ，其类 \(\widetilde{f}_1\) 与 \(\widetilde{f}_2\) 在 \(\mathrm{L}^2 (\mathrm{X},\mu)\) 中。由于 \(f(m_g) = m_{f\circ g}\) （命题 4, c)），\((\widetilde{f}_1,\widetilde{f}_2)\) 关于 \(f(m_g)\) 的谱测度是像测度 \((f\circ g)(\overline{f}_1f_2\cdot \mu)\) （IV, p. 269 的引理 4, b)）。这个测度等于像测度 \(f(g(\overline{f}_1f_2\cdot \mu))\) （INT, V, p. 72, § 6, 第 4 目, 命题 4, a)），由此即得断言（IV, p. 269 的引理 4, b)）。

推论 4. --- 设 \(g \in \mathcal{L}_{\mathrm{u}}(\mathrm{Sp}(f(u)))\) 。我们有 \(g(f(u)) = (g \circ f)(u)\) 。

我们有 \(g \circ f \in \mathcal{L}_{\mathrm{u}}(\mathrm{Sp}(u))\) （IV, p. 184 的引理 5）。对 E 中所有 \(x\) 与 \(y\) ，用 \(\mu_{x,y}^{\prime}\) 表示 \((x,y)\) 关于 \(f(u)\) 的谱测度。由前面的推论以及 INT, V, p. 71, § 6, 第 2 目, 定理 1，我们有 \(g \in \mathcal{L}^{2}(\mathrm{Sp}(f(u)), \mu_{x,y}^{\prime})\) 当且仅当 \(g \circ f \in \mathcal{L}^{2}(\mathrm{Sp}(u), \mu_{x,y})\) ，因此 \(g(f(u))\) 的定义域等于 \((g \circ f)(u)\) 的定义域。于是，对所有 \(x \in \mathrm{E}\) 及 \(y \in \operatorname{dom}(g(f(u)))\) ，有公式

\[
\begin{array}{c} \langle x \mid g (f (u)) y \rangle = \int_ {\mathrm{Sp} (f (u))} g \mu_ {x, y} ^ {\prime} \\ = \int_ {\mathrm{Sp} (f (u))} g f (\mu_ {x, y}) = \langle x \mid (g \circ f) y \rangle \end{array}
\]

（loc. cit.），由此 \(g(f(u)) = (g \circ f)(u)\) 。

命题 5. --- 设 u 是希尔伯特空间 E 上的正规部分算子。

\begin{enumerate}
\def\labelenumi{\alph{enumi})}
\item
  \(If f \in \mathcal{L}_{\mathrm{u}}^{\infty}(\mathrm{Sp}(u)),\) 则 \(f(u) \in \mathcal{L}(\mathrm{E})\) ；
\item
  \(If f \in \mathcal{L}_{\mathrm{u}}^{\infty}(\mathrm{Sp}(u))\) 且 \(g \in \mathcal{L}_{\mathrm{u}}(\mathrm{Sp}(u))\) ，则 \(f(u) \circ g(u) \subset (fg)(u)\) ；
\item
  映射 \(f \mapsto f(u)\) 从 \(\mathcal{L}_{\mathrm{u}}^{\infty}(\mathrm{Sp}(u))\) 到 \(\mathcal{L}(\mathrm{E})\) 中是星代数的连续单位态射。特别地，有 \(\|f(u)\|\leqslant\|f\|_{\infty}\) ，其中 \(f\in\mathcal{L}_{\mathrm{u}}^{\infty}(\mathrm{Sp}(u))\) ；
\item
  若 \(u \in \mathcal{L}(\mathrm{E})\) ，则对所有 \(f \in \mathcal{L}_{\mathrm{u}}^{\infty}(\mathrm{Sp}(u))\) ，自同态 \(f(u)\) 属于 \(u\) 在 \(\mathcal{L}(\mathrm{E})\) 中的双交换子。
\end{enumerate}

设 \(f \in \mathcal{L}_{\mathrm{u}}^{\infty}(\mathrm{Sp}(u))\) 。由 p. 272 的公式 (8)，我们有 \(\mathrm{D}_f = \mathrm{E}\) ，且 \(f(u)\) 是 \(\mathrm{E}\) 的连续自同态，由此得断言 \(a\) )。

设 \(g \in \mathcal{L}_{\mathrm{u}}(\mathrm{Sp}(u))\) 。设 \(y \in \mathrm{D}_g\) 。我们有 \(y \in \mathrm{D}_{fg}\) ，并且对所有 \(x \in \mathrm{E}\) ，由 p. 271 的公式 (7) 可得 \(\langle \overline{f}(u)x | g(u)y \rangle = \langle x | (fg)(u)y \rangle\) ，由此 \(f(u)(g(u)y) = (fg)(u)(y)\) ，这就证明了 \(b\) 。

由 \(a\) ) 与 \(b\) )，映射 \(f \mapsto f(u)\) 从 \(\mathcal{L}_{\mathrm{u}}^{\infty}(\mathrm{Sp}(u))\) 到 \(\mathcal{L}(\mathrm{E})\) 中，是对合代数的单位态射；因而它是范数 \(\leqslant 1\) 的连续态射（I, p. 104, 命题 2），由此得断言 \(c\) )。

假设 \(u\) 是 E 的自同态。设 \(v \in \mathcal{L}(\mathrm{E})\) 与 \(u\) 交换。于是，由连续函数演算的性质（I, p. 110, 注），我们有 \(v \circ f(u) = f(u) \circ v\) ，这里 \(f \in \mathcal{C}(\mathrm{Sp}(u))\) 。设 \(x\) 与 \(y\) 属于 E。公式

\[
\langle x \mid (v \circ f (u)) y \rangle = \langle x \mid (f (u) \circ v) y \rangle ,\tag{9}
\]

对每个函数 \(f \in \mathcal{C}(\mathrm{Sp}(u))\) 成立，它表示 \((v^{*}(x), y)\) 与 \((x, v(y))\) 关于 u 的谱测度相等。由 p. 271 的公式 (6)，这个等式蕴含公式 (9) 对所有 \(f \in \mathcal{L}_{\mathrm{u}}^{\infty}(\mathrm{Sp}(u))\) 成立。因此我们有 \(v \circ f(u) = f(u) \circ v\) 。

推论. --- 设 \(f\) 与 \(g\) 属于 \(\mathcal{L}_{\mathrm{u}}(\mathrm{Sp}(u))\) ，并设 \((x,y)\in\mathrm{D}_{f}\times\mathrm{D}_{g}\) 。\((f(u)x,g(u)y)\) 关于 \(u\) 的谱测度是测度 \(\overline{f}g\cdot\mu_{x,y}\) ，其中 \(\mu_{x,y}\) 是 \((x,y)\) 关于 \(u\) 的谱测度。

用 \(\nu\) 表示 \((f(u)x,g(u)y)\) 关于 \(u\) 的谱测度。对每个函数 \(\varphi \in \mathcal{H}(\mathrm{Sp}(u))\) ，我们有

\[
\begin{array}{r l} & {\int_ {\mathrm{Sp} (u)} \varphi \nu = \langle f (u) x | \varphi (u) (g (u) y) \rangle} \\ & {\qquad = \langle f (u) x | (\varphi g) (u) y \rangle = \int_ {\mathrm{Sp} (u)} \overline {{f}} \varphi g \mu_ {x, y},} \end{array}
\]

（命题 5, b)）由此 \(\nu = \overline{f} g \cdot \mu_{x,y}\) 。

命题 6. --- 设 u 是希尔伯特空间 E 上的正规部分算子。设 \((f_n)_{n \in \mathbf{N}}\) 是 \(\mathcal{L}_{\mathrm{u}}(\mathrm{Sp}(u))\) 中的序列，它逐点收敛于 \(f \in \mathcal{L}_{\mathrm{u}}(\mathrm{Sp}(u))\) ，且存在 \(g \in \mathcal{L}_{\mathrm{u}}(\mathrm{Sp}(u))\) 满足 \(|f_n| \leqslant g\) ，对所有 \(n \in \mathbf{N}\) 。这时我们有 \(\operatorname{dom}(g(u)) \subset \operatorname{dom}(f_n(u))\) ，对所有 \(n \in \mathbf{N}\) ，以及 \(\operatorname{dom}(g(u)) \subset \operatorname{dom}(f(u))\) 。此外，对 \(g(u)\) 的定义域的每个元素 y，我们有

\[
f (u) (y) = \lim _ {n \to + \infty} f _ {n} (u) (y).
\]

特别地，若 \(f_{n} \in \mathcal{L}_{\mathrm{u}}^{\infty}(\mathrm{Sp}(u))\) 对所有 \(n \in \mathbf{N}\) 成立，且函数 \(f_{n}\) 一致有界，则 \(f_{n}(u)\) 收敛于 \(f(u)\) ，这里的空间 \(\mathcal{L}(\mathrm{E})\) 赋以逐点收敛拓扑。

用 \(\mu_{x,y}\) （相应地，\(\mu_{x}\) ）表示 \((x,y)\) （相应地，x）关于 u 的谱测度。

设 \(y \in \mathrm{dom}(g(u))\) ，于是 \(g \in \mathcal{L}^2(\mathrm{Sp}(u), \mu_y)\) 。条件 \(|f_n| \leqslant g\) 蕴含 \(f_n \in \mathcal{L}^2(\mathrm{Sp}(u), \mu_y)\) ，因此 \(y \in \mathrm{dom}(f_n(u))\) 。

由于 \((f_{n})\) 逐点收敛于 f 且 \(|f_{n}| \leqslant g\) ，由勒贝格定理（INT, IV, p. 137, § 3, 第 7 目, 定理 6），我们有 \(f \in \mathcal{L}^{2}(\mathrm{Sp}(u), \mu_{y})\) ，因此 \(y \in \mathrm{dom}(f(u))\) 。此外，序列 \((f_{n})\) 在 \(\mathcal{L}^{2}(\mathrm{Sp}(u), \mu_{y})\) 中收敛于 f，于是 \(f_{n}(u)y\) 的范数收敛于 \(f(u)y\) 的范数。

设 \(x \in \mathrm{E}\) 。函数 \(f\) 与 \(g\) ，以及函数 \(f_n\) （对所有 \(n \in \mathbf{N}\) ），都属于 \(\mathcal{L}^1(\mathrm{Sp}(u), \mu_{x,y})\) （命题 3）。由勒贝格定理（INT, IV, loc. cit.），序列 \((f_n)\) 收敛于 \(f\) （在 \(\mathcal{L}^1(\mathrm{Sp}(u), \mu_{x,y})\) 中），由此

\[
\langle x \mid f _ {n} (u) y \rangle = \int_ {\operatorname{Sp} (u)} f _ {n} \mu_ {x, y} \rightarrow \int_ {\operatorname{Sp} (u)} f \mu_ {x, y} = \langle x \mid f (u) y \rangle .
\]

由此得出 \(f_{n}(u)(y)\) 收敛于 \(f(u)(y)\) （EVT, V, p. 17, 命题 10）。

命题 7. --- 设 X 是局部紧拓扑空间，ν 是 X 上的测度。设 g: C × X → C 是具有紧支撑的连续函数。对 z ∈ C，置

\[
h (z) = \int_ {\mathrm{X}} g (z, x) d \nu (x).
\]

\begin{enumerate}
\def\labelenumi{\alph{enumi})}
\item
  从 C 到 C 的映射 h 是连续且有界的；
\item
  从 X 到 \(\mathcal{L}(\mathrm{E})\) 的、由 \(x\mapsto g(u,x)\) 定义的映射是 \(\nu\)-可积的，并且我们有
\end{enumerate}

\[
h (u) = \int_ {\mathrm{X}} g (u, x) d \nu (x).
\]

函数 h 是有界的，因为 g 是具有紧支撑的连续函数；且由 INT, IV, p. 144, § 4, 第 3 目, 推论 1，它是连续的。由于 g 连续且具有紧支撑，映射 \(x \mapsto g(u, x)\) 从 X 到 \(\mathcal{L}(\mathrm{E})\) 中连续。这个映射具有紧支撑，故在 X 上关于 \(\nu\) 有界且可积。我们记

\[
v = \int_ {\mathrm{X}} g (u, x) d \nu (x) \in \mathscr {L} (\mathrm{E}).
\]

设 y 与 z 是 E 的元素，并设 \(\mu\) 是 \((y, z)\) 关于 u 的谱测度。我们有

\[
\langle y \mid v (z) \rangle = \int_ {\mathrm{X}} \langle y \mid g (u, x) z \rangle d \nu (x) = \int_ {\mathrm{X}} \left(\int_ {\operatorname{Sp} (u)} g (\lambda , x) d \mu (\lambda)\right) d \nu (x)
\]

（p. 271 的公式 (6)）。由于 \(g \in \mathcal{K}(\mathbf{C} \times \mathbf{X})\) ，由此可得

\[
\langle y \mid v (z) \rangle = \int_ {\operatorname{Sp} (u)} \Bigl (\int_ {\mathrm{X}} g (\lambda , x) d \nu (x) \Bigr) d \mu (\lambda) = \langle y \mid h (u) z \rangle
\]

依据 INT, III, p. 84, § 4, 第 1 目, 定理 2 以及 p. 271 的公式 (6)。这就证明了 \(v = h(u)\) ，此即所要证者。

